\setcounter{section}{44}
\section{Esquema Nacional de Interoperabilidad}

Nombre completo: El Esquema Nacional de Interoperabilidad. Dimensiones de la interoperabilidad.

Apuntes de ASTIC de 2023.

\localtableofcontents
\newcommand{\ISA}{ISA\textsuperscript{2}}

\subsection{Esquema Nacional de Interoperabilidad} %7
    \subsubsection{Introducción} %7
        El Esquema Nacional de Interoperabilidad es el marco que rige el intercambio y conservación de la información en las administraciones públicas. En este documento colaboran todas las Administraciones Públicas a través de órganos Colegiados con competencia en administración digital, y la opinión de asociaciones de la industria. Formando el \textbf{Real Decreto 4/2010}.

        Tiene una relación estrecha con la seguridad y debe abordarse desde una perspectiva integral y coordinada. El impulso de la interoperabilidad se vió reflejado en la ley 40/2015 de régimen jurídico del sector público. Esta determina que las Administraciones Públicas deben relacionarse entre sí a través de medios electrónicos que aseguren la interoperabilidad. Establece las condiciones en los artículos 155 a 158.

        La interoperabilidad es una cualidad integral en la que intervienen múltiples elementos que deben coordinarse en su conjunto.

    \subsubsection{Normativa y marco jurídico aplicable} %8
        Se deben tener en cuenta las siguientes normas en el entorno del Esquema Nacional de Interoperabilidad:
        \begin{description}
            \item[RD 4/2020] por el que se regula el Esquema Nacional de Interoperabilidad
            \item[RD 203/2021] Por el que se aprueba el Reglamento de actuacióny funcionamiento del sector público por medios electrónicos
            \item[L 39/2015] Procedimiento administrativo común
            \item[L 40/2015] Régimen Jurídico del sector público
            \item[RUE 2016/679] RGPD
            \item[LO 3/2018] de Protección de datos
            \item[RD 311/2022] Esquema Nacional de Seguridad
            \item[RD 1112/2018] Accesibilidad de los sitios web
            \item[RD 1494/2027] condiciones básicas para el acceso de las personas con discapacidad
            \item[L 19/2013] LTBG
            \item[37/2007] reutilización de información del sector público
            \item[RD 1708/2011] Sistema Español de Archivos
        \end{description}

    \subsubsection{Objetivos} %9
        Con el Esquema Nacional de Interoperabilidad se persiguen los siguientes objetivos:
        \begin{itemize}
            \item Establecer y detallar los criterios que deben seguir las administraciones públicas para la toma de decisiones tecnológicas que garanticen la interoperabilidad y que eviten la discriminación a los ciudadanos por razón de su elección tecnológica.
            \item Contribuir a la interoperabilidad en el uso de medios electrónicos que permitan a los ciudadanos y administraciones el ejercicio de sus derechos y cumplimiento de deberes públicos en materia de interoperabilidad.
            \item Proporcionar los elementos comunes que han de facilitar la actuación de las Administraciones Públicas en materia de Interoperabilidad
            \item Facilitar la implantación de las políticas de seguridad al contribuir a  un escenario de mayor racionalidad técnica y de economía de escala.
        \end{itemize}

    \subsubsection{Ámbito de aplicación} %9
        Se establece en el Art. 2 de las leyes 39/2015 y 40/2015 que determina su aplicación en la AGE, Administración de las Comunidades Autónomas, Entidades integradas en la Administración Local y el Sector público Institucional. Este último integrado por organismos públicos y entidades de derecho público vinculadas o dependientes de las AAPP, Entidades de derecho privado vinculadas o dependientes de las AAPP, sujetas a lo dispuesto en la ley y las Universidades públicas (que tendrán su propia norma supletoria a lo previsto en la ley).

    \subsubsection{Principios básicos de la interoperabilidad} %10
        Los principios básicos egún el Art. 4, y definidos en los ARt. 5, 6 y 7 son:
        \begin{itemize}
            \item Interoperabilidad como cualidad integral presente desde la concepción y durante toda su vida útil
            \item Carácter multidimensional de la interoperabilidad
            \item Soluciones multilaterales
        \end{itemize}

    \subsubsection{Interoperabilidad en la unión europea} %11
        La cooperación es una cualidad necesaria de la Unión Europea, con las Administraciones de los diferentes estados debiendo cooperar para resolver los problemas de los ciudadanos y de las instituciones comunitarias en cuestiones de políticas comunes.

        La interoperabilidad es necesaria para ofrecer servicios públicos digitales transfronterizos, que concluye en el principio de sólo una vez (once only). Esto se ve culminado en el Marco Europeo de Interoperabilidad, la Declaración Ministerial de Tallin, la Decisión \ISA y los Planes de Acción de administración electrónica.

        La \textbf{Declaración Ministerial de administración electrónica de Tallín} adoptada en 2017, refuerza el compromiso con los principios del Plan de Acción de Administración electrónica 2016-2020. Incluye líneas de acción relativas al principio de solo una vez y a la interoperabilidad por defecto. Busca promover la reutilización de datos en los registros primarios y de las soluciones del ámbito comunitario, CEF Telecom e \ISA y de las administraciones de los estados miembros. 

        \paragraph{Marco Europeo de Interoperabilidad} %11
            La Comisión Europea identificó la necesidad de interoperabilidad entre administraciones públicas en 1999. Ha respaldado programas para el desarrollo de soluciones de interoperabilidad. En 2010 adoptó una comunicación que adjuntaba la Estrategia Europea de Interoperabilidad y el Marco Europeo de Interoperabilidad, que sirve de base para la mayoría de marcos nacionales de interoperabilidad.

            En 2017 se publicó una nueva comunicación para adaptarse a los probramas iniciados y a las tendencias y desarrollos tecnológicos. Los programas \textbf{Soluciones de Interoperabilidad para las Administraciones} ISA (2010-2015) e \ISA (2016-2020) fueron los principales instrumentos para los que se han aplicado la Estrategia Europea de Interoperabilidad y el Marco Europeo de Interoperabilidad actuales.
            
            Se establece la Interoperabilidad como:

            \noindent\fbox{\parbox{\textwidth}{
                A efectos del EIF, la interoperabilidad es la capacidad de que las organizaciones interactúen con vistas a alcanzar objetivos comunes que sean mutuamente beneficiosos y que hayan sido acordados previa y conjuntamente, recurriendo a la puesta en común de información y conocimientos entre las organizaciones, a través de los procesos empresariales a los que apoyan, mediante el intercambio de datos entre sus sistemas TIC respectivos.
            }}

            La finalidad del Marco Europeo de Interoperabilidad queda fijada en:
            \begin{itemize}
                \item Orientar a las Administraciones públicas europeas en sus esfuerzos por diseñar y prestar sin discontinuidad a otras administraciones públicas, ciudadanos o empresas servicios públicos europeos que sean, en la medida de lo posible, en principio digitales, transfronterizos y abiertos.
                \item Orientar a las administraciones públicas en la elaboración y actualización de los Marcos Nacionales de Interoperatividad, o políticas, estrategias y orientaciones nacionales para promover la Interoperatividad.
                \item Contribuir a la creación de un mercado único digital para la prestación de servicios públicos europeos.
            \end{itemize}
            La comuicación reconoce la dificultad de adaptar la situación y especificaciones nacionales al contexto europeo.

            El Marco Europeo de Interoperabilidad presenta los siguientes doce principios:
            \begin{description}
                \item[Subsidiariedad y proporcionalidad] Tomar las decisiones cerca del ciudadano y limitar la acción de la UE a lo necesario
                \item[Apertura] en Datos, especificaciones y Software
                \item[Transparencia]
                \item[Posibilidad de reutilización]
                \item[Neutralidad tecnológica y portabilidad de los datos]
                \item[Primacía del usuario]
                \item[Inclusión y accesibilidad]
                \item[Seguridad e intimidad]
                \item[Multilingüismo]
                \item[Simplificación administrativa]
                \item[Conservación de la información]
                \item[Evaluación de efectividad y Eficiencia]
            \end{description}
            Estos puntos se explican en una serie de 47 recomendaciones con mas detalle.
            
        \paragraph{Programa \ISA} %14
            Este programa vigente entre 2016 y 2020 continúa la labor de programas previos, facilitando la interacción electrónica transfronteriza, tanto entre las Administraciones Públicas, como con las empresas y ciudadanos. Destacando el mantenimiento de la estrategia, marco y arquitectura europea de interoperabilidad, la promoción de soluciones de interoperabilidad y su reutilización, mecanismos para cuantificar ahorros de costes por interoperabilidad y detección de lagunas legislativas y evaluación de implicaciones TIC en la legislación comunitaria.

        \paragraph{Mecanismo Conectar Europa (CEF)} %14
            Connect Europe Facility, se lanzó en 2013 para finalizar el despliegue de infraestructuras pan-europeas de interconexión de estados miembros en temas de energía, transporte y telecomunicaciones (CEF Telecom). Este último financia infraestructuras de servicios digitales para la interacción entre autoridades públicas, ciudadanos y empresas a través de plataformas de interacción normalizas.

            Estas infraestructuras consideradas componentes elementales son: identificación y autenticación electrónica, entrega electrónica de documentos, traducción automática, archivo electrónico y ciberseguridad y facturación electrónica

        \paragraph{Arquitectura Europea de Interoperabilidad (EIRA)} %14
            La Comisión Europea proponeun metamodelo basado en el Marco Europeo de Interoperabilidad: Arquitectura Europea de Interoperabilidad (EIRA).

            Tiene como objetivo asegurar interoperabilidad técnica, sintáctica, semántica, organizativa y legal entre datos, procesos y servicios digitales.

            EIRA define los bloques de construcción arquitectónicos necesarios para construir sistemas de administración electrónica interoperables y proporciona terminología común.

            Su versión 6 fue publicada en Mayo de 2023.

        \paragraph{Observatorio de los Marcos Nacionales de Interoperabilidad (NIFO)} %15
            En el contexto del programa \ISA se estableció el NIFO para medir los progresos y supervisar el estado de la interoperabilidad y el progreso de los países. La Comisión Europea a través del NIFO publica los siguientes informes:
            \begin{description}
                \item[State of play report] proporciona una panorámica de los progresos conseguidos emn Administración Digital e Interoperabilidad, así como las medidas implementadas por la Comisión para facilitar a los países alcanzar sus objetivos en estos ámbitos.
                \item[Fichas NIFO] Analizan el alineamiento de los estados miembros con el marco europeo de interoperabilidad
                \item[Informe LIFO] visión general de la información recopilada sobre interoperabilidad.
            \end{description}

            El mecanismo de Seguimiento del EIF está organizado en los pilares de Principios, Niveles y Modelo, agrupando 47 recomendaciones y 68 KPIs.

            El informe de situación se usa con la toma de indicadores primarios recopilados en una encuesta difundida a los puntos de contacto nacionales, e indicadores secundarios recopulados de fuentes externas. Midiendo la implementación de las 47 recomendaciones y 68 KPIs del EIF.

            El último informe de situación es de 2022, las últimas fichas NIFO son de 2023 y el último informe LIFO es de 2020. Ocupando España el octavo lugar.

        \paragraph{Resumen de posicionamiento de España en 2022} %16
            España lidera la media de resultados en los campos analizados %Esta frase es de que quien no se contenta es porque no quiere.
            En la mayoría de los indicadores España se situa por encima de la media europea.

            Destaca la adopción de España de la Carta de Derechos Digitales, consagrando principios y derechos que orientarán los proyectos normativos y el desarrollo de políticas públicas para garantizar la protección de derechos individuales y colectivos en los escenarios digitales.

            También se destaca el principio de sólo una vez, promovido por las leyes 39/2015 y 40/2015. Destacando el uso de la Plataforma de Intermediación de Datos.

        \paragraph{JOINUP: Plataforma colaborativa creada por la Comisión Europea} %17
            Creada por la Comisión Europea para proporcionar un lugar común para las administraciones públicas, empresas y ciudadanos para compartir y reutilizar soluciones y buenas prácticas, y facilitar la comunicación con Europa.

        \paragraph{SEMIC: Semantic Interoperability Comunity} %17
            La interoperabilidad semántica es vital para la interoperabilidad entre administraciones. Por ello se promueven políticas de metadatos comunes, transparentes y bien documentadas.

            SEMIC permite colaborar y compartir herramientas, estándares y buenas prácticas. Ofrece la publicación de especificaciones, desarrolla soluciones, provee herramientas, formación y concienciación. Lo hace a través de:
            \begin{itemize}
                \item Registros de base
                \item Catálogo de servicios
                \item Espacio de datos.
            \end{itemize}

    \subsubsection{Infraestructura y servicios comunes} %18
        En el Art. 12 se establece que las Administraciones públicas enlazarán aquellas infraestructura y servicios que puedan implantar con las infraestructuras comunes que porporcione la AGE.

        El \textbf{catalogo de servicios comunes, infraestructuras y otras soluciones}, ofrecido por la Secretaría General de la Administración Digital está estructurado en los siguientes apartados:
        \begin{itemize}
            \item Administración digital y servicios al ciudadano
            \item Gestión interna
            \item Infraestructuras
            \item Regulación, guías e informes relevantes
        \end{itemize}

        Destacan las infraestructuras comunes y servicios compartidos, salvo las excepciones del marco regulador en relación con la Seguridad Social, la informática presupuestaria de la Intervención General del Estado y la Agencia Estatal de la Administración Tributaria, Defensa y las Fuerzas y Cuerpos de Seguridad del Estado.

        Los beneficios esperados son:
        \begin{itemize}
            \item Ahorro de costes potencial
            \item Eliminación de duplicidades
            \item Mejora de servicios prestados a los ciudadanos
            \item Mejora de calidad o funcionalidad en procesos internos
            \item Adquisición de capacidades o recursos
            \item Liberación de recursos humanos
            \item Liberación de recursos económicos para iniciativas de innovación o transformación digital
        \end{itemize}

        Los servicios de caracter obligatorio son los siguientes:
        \begin{itemize}
            \item Servicio unificado de telecomunicaciones
            \item SErvicio de seguridad gestionada
            \item Servicio de alojamiento de infraestructuras TIC
            \item Servicio de nube híbrida (Nube SARA)
            \item Correo electrónico
            \item Servicio multicanal de atención al ciudadano
            \item Servicio de gestión del registro
            \item Servicio de gestión de notificaciones
            \item Servicio de gestión de nómina
            \item Servicio integrado de gestión de personal
            \item Servicio común de gestión economico-presupuestaria
            \item Servicio común de generación y validación de firmas electrónicas
            \item Servicio de gestión de expediente y documento electrónico
            \item Servicio de gestión de archivo electrónico
        \end{itemize}

    \subsubsection{Comunicaciones de las Administraciones Públicas} %19
        El capítulo VII del título \textit{Comunicaciones de las Administraciones públicas} establece que se usará preferentemente la Red de comunicaciones de las Administraciones públicas españolas para comunicarse entre sí. Para ello conectarán a la misma sus redes o sus nodos de interoperabilidad, así como la conexión con las redes de las instituciones de la Unión Europea y otros estados miembros. En España esto lo realiza la red SARA.

        Se debe cumplir lo establecido en la NTI de Requisitos de conexión y la guía de aplicación a la red de comuncaciones de las Administraciones Públicas. Estas aplicarán el plan de direccionamiento e interconexión de redes en la administración.

        Los sistemas y aplicaciones implicados se sincronizan con la hora oficial, que garantiza la certidumbre de plazos establecidos. Esto se realiza con el Real Instituto y Observatorio de la Armada 

    \subsubsection{Dimensiones de la interoperabilidad} %20
        Son objeto del tema 046

    \subsubsection{Normas Técnicas} %20
        Son objeto del tema 046

    \subsubsection{Instrumentos para la interoperabilidad} %20
        Tras el RD 203/2021 quedan establecidos los siguientes instrumentos para la interoperabilidad:
        \begin{itemize}
            \item Sistema de Información Administrativa
            \item Centro de Transferencia Tecnológica
            \item Directorio Común de Unidades Orgánicas y Oficinas de las Administraciones Públicas
            \item Centro de interoperabilidad semántica de la administración
        \end{itemize}

        \paragraph{Sistema de Información Administrativa}%20
            Es el inventario de procedimientos administrativos, servcicios prestados y otras actuaciones administrativas que generan documentación pública. Contiene información clasificada por funciones e indica su nivel de informatización y las interfaces para favorecer la interacción o integración con procesos.

        \paragraph{Directorio Común de Unidades Orgánicas y Oficinas de las Administraciones Públicas} %20
            Permite la sincronización de sistemas de información, inventariando, codificando y evolucionando las unidades orgánicas y oficinas según modalidades e integración. Busca la flexibilidad en el consumo y en la provisión de información

        \paragraph{Centro de transferencia de tecnología} %21
            Se describe el CTT como el directrio de aplicaciones par su libre reutilización. Se registra la aplicación, descripción, uso, características, licencia, estándares usados y estado del desarrollo.

            Según el modelo de colaboración, existen modalidades de descarga de los activos. Para la descarga es importante confirmar la idoneidad del uso. Están presentes en Github como \href{https://github.com/ctt-gob-es}{Forja-CTT}. Incluye clientes de firma, o de acceso a tarjetas inteligentes.

            Tiene un buscador propio para soluciones clasificado. Las aplicaciones deben estar licenciadas. La licencia debe ser EUPL o compatible.

        \paragraph{Centro de Interoperabilidad semántica de la Administración} %23
            El CISE es el instrumento definido para publicar modelos de datos de los elementos de interoperabilidad que permiten intercambiar información. Sus principales objetivos son:
            \begin{itemize}
                \item Identificar modelos de datos que tienen carácter de comunes
                \item Facilitar la comprensión semántica de los servicios de intercambio de datos de las Administraciones
                \item Maximizar la reutilización de activos semánticos en su construcción
                \item Proporcionar un repositorio de modelos de datos de intercambio de información, debe permitir recopilarlos, clasificarlos, publicarlos, localizarlos y difundirlos.
                \item Facilitar la interoperabilidad publicando modelos de datos de los servicios de intercambio mediante esquemas XML (XSD), sus codificaciones y documentación asociada.
            \end{itemize}

            Se denomina \textbf{modelo de datos} al conjunto de definiciones, interrelaciones, reglas y convenciones que permiten describir los datos para su intercambio.

            Un modelo de datos se compone de archivos semánticos, piezas reutilizables que habilitan la estandarización de conceptos. Son los esquemas xml, los que permiten definir cada archivo semántico en una serie de campos complejos y simples.

            \begin{lstlisting}[language=XML,gobble=16]
                <?xml version="1.0" encoding="UTF-8"?> 
                <xsd:schema xmlns:xsd="http://www.w3.org/2001/XMLSchema">
                    <xsd:element name="Libro">
                        <xsd:complexType>
                            <xsd:sequence>
                                <xsd:element name="Titulo" type="xsd:string"/>
                                <xsd:element name="Autores" type="xsd:string" maxOccurs="10"/>
                                <xsd:element name="Editorial" type="xsd:string"/>
                            </xsd:sequence>
                            <xsd:attribute name="precio" type="xsd:double"/>
                        </xsd:complexType>
                </xsd:element>
                </xsd:schema>
            \end{lstlisting}
            
    \subsubsection{Reutilización y transferencia de tecnología} %24
        La reutilización es objeto de los artículos 16 y 17.
        
        En la ley 40/2015 del régimen jurídico del sector público, ens u ARt. 157 se establece que:
        \begin{itemize}
            \item Las administraciones pondrán a disposición de cualquiera de ellas que lo solicite las aplicaciones desarrolladas y de cuyos derechos de popiedad sean titulares. Las administraciones cedentes y cesionarias podrán acordar la repoercusión del coste. Estas aplicaciones podrán ser declaradas como de fuente abierta
            \item Las Administraciones públicas con carácter previo a la adquisición, desarrollo, o mantenimiento a lo largo de todo el ciclo de vida, deberá consultar en el directorio general de aplicaciones, dependiente de la AGE si existen soliuciones disponibles para su reutilización.
            \item (Art. 158) Las administraciones Públicas mantendrán directorios actualizados de aplicaciones para su libre reutilizacións, que deberán ser plenamente interoperables para garantizar su compatibilidad informática e interconexión.
        \end{itemize}

        Se considera activo reutilizable cualquier tipo de material electrónico que pueda ser empleado de forma recurrente. Algunos ejemplos son programas de ordenador, guías, manuales, interfaces, modelos de datos, algoritmos, patrones de diseño, arquitecturas, especificaciones de requerimientos de diseño y prueba, guías de uso, información del sector público, vídeos, audios, materiales divulgativos\dots

        El Art. 16 describe las condiciones de licenciamiento de las aplicaciones informáticas, documentación y cualquier otro objeto de información cuya titularidad sea una Administración Pública y que:
        \begin{itemize}
            \item el fin es el aprovechamiento de los recursos públicos
            \item protección contra la apropiación total o parcial por parte de terceros
            \item exención de responsabilidad del cedente por el mal uso de los cesionarios
            \item no obligación de asistencia técnica o mantenimiento del cedente
            \item ausencia de responsabilidad del cedente en caso de errores o mal funcionamiento
            % \item Lo que parece una licencia permisiva del copón de la baraja
        \end{itemize}

        El licenciamiento es por defecto, sin contraprestación y sin convenio. solo se podrá acordar la repercusión parcial del coste de adquisición en los casos en los que repoercuta directamente en el incremento de funciones, adaptaciones o implique el suministro de servicios de asistencia o soporte para su reutilización en el organismo cesionario.

        Las licencias deberán cumplir con los 4 requisitos de SW libre de Stallman:
        \begin{itemize}
            \item Ejecutarse para cualquier propósito
            \item Conocer su código fuente
            \item Modificarse o mejorarse
            \item Redistribuir con o sin cambios, manteniendo estas garantías
        \end{itemize}

        Se procura para este fin la Licencia Pública de la Unión Europea (Una licencia permisiva compatible con otras licencias similares). Esta licencia se aplica a software, especificaciones, documentación técnica y estándares.

        La administración contratante podrá incluir en los pliegos los siguientes aspectos: Incluir los derechos completos de propiedad intelectual y recibir un producto que pueda ofrecer para su posterior reutilización y declarar como fuente abierta.

        Se establece en el Art.17 el mantenimiento del Directorio General de Aplicaciones para su libre reutilización, a través del Centro de Transferencia de Tecnología. Las administraciones públicas conectarán los directorios de aplicaciones para su libre reutilización entre sí y con la Unión Europea.

    \subsubsection{Reutilización de recursos de la información} %26
        Uso de datos generados y custodiados por organismos del sector público con fines comerciales o no. Los beneficios son: el potencial económico de desarrollo de nuevos productos: fomento económico y labotal e incremento de la transparencia administrativa.

        La iniciativa Aporta, creada en el 2009, tiene tres objetivos:
        \begin{itemize}
            \item Fomento de una cultura favorable a los datos públicos abiertos
            \item Facilitar la apertura de las administraciones
            \item Impulsar el mercado de la reutilización de información pública
        \end{itemize}

        El portal \url{datos.gob.es} tiene caracter nacional y gestiona el Catalogo público. También es el canal de novedades, formación y concienciación. Mas información en el tema siguiente.










        
%     \subsubsection{Firma electrónica y certificados} %27
%     \subsubsection{Recuperación y Conservación del Documento Electrónico} %28
%     \subsubsection{Conformidad con el Esquema Nacional de Interoperabilidad} %30
%     \subsubsection{Conclusiones} %35

% \subsection{Dimensiones de la interoperabilidad} %38
%     \subsubsection{Introducción} %38
%     \subsubsection{Normativa o marco jurídico aplicable} %38
%     \subsubsection{Niveles de Interoperabilidad en Europa} %38
%         \paragraph{Interoperabilidad Jurídica} %38
%         \paragraph{Interoperabilidad Organizativa} %39
%         \paragraph{Interoperabilidad Semántica} %39
%         \paragraph{Interoperabilidad Técnica} %39
%     \subsubsection{Dimensiones} %40
%         \paragraph{Dimensión Legal} %40
%         \paragraph{Dimensión Temporal} %40
%         \paragraph{Dimensión Organizativa} %41
%         \paragraph{Dimensión Semántica} %43
%         \paragraph{Dimensión Técnica} %43